\documentclass[11pt]{article}
\usepackage[margin=1in]{geometry}
\usepackage{amsmath,amssymb,amsthm}
\usepackage{graphicx}
\usepackage[hidelinks]{hyperref}

\newtheorem{theorem}{Theorem}
\newtheorem{remark}{Remark}

\newcommand{\BPBp}{\mathrm{BPBp}}
\newcommand{\K}{\mathbb K}
\newcommand{\eps}{\varepsilon}
\newcommand{\ot}{\widehat{\otimes}_{\varepsilon}}

\title{Every nonzero injective tensor factor descends the Bishop-Phelps-Bollobas property}
\author{Run \texttt{fa\_banach\_001}, lane 5}
\date{2026-06-19}

\begin{document}
\maketitle

\section*{Status}

This packet gives a full qualitative answer to the tensor-factor descent
question following Proposition 2.8 in Aron, Choi, Kim, Lee, and
Martin~\cite{ACKLM}.  The source proves the result for $C(K)$ and asks what
other spaces have the same descent property.  The theorem below shows that,
as a yes/no BPBp statement, every nonzero Banach tensor factor $E$ works.

There is a modulus nuance.  The source's $C(K)$ argument preserves the same
modulus $\eta(\eps)$.  The all-$E$ theorem below gives the modulus
$\eps\mapsto \eta(\eps/2)$.  A same-modulus refinement still holds under the
stronger state-norming-unit, or spear-vector, hypothesis discussed at the end.

\section*{Source question}

The relevant passage appears on page 11 of the source paper, immediately after
Proposition 2.8.

\begin{center}
\includegraphics[width=\textwidth]{figures/open_problem_crop.png}
\end{center}

Transcription of the part used here: Proposition 2.8 may be restated as saying
that if $(X,C(K)\ot Y)$ has the BPBp with $\eta(\eps)$, then so does $(X,Y)$.
The authors write that they do not know what other spaces, besides $C(K)$,
have this property.

\section*{Definitions}

All Banach spaces are over $\K=\mathbb R$ or $\mathbb C$.  A pair $(X,Y)$ has
the Bishop-Phelps-Bollobas property with function $\eta$ if, whenever
$T\in S_{\mathcal L(X,Y)}$ and $x_0\in S_X$ satisfy
\[
 \|Tx_0\|>1-\eta(\eps),
\]
there are $S\in \mathcal L(X,Y)$ and $x_1\in S_X$ such that
\[
 \|S\|=\|Sx_1\|=1,\qquad \|S-T\|<\eps,\qquad
 \|x_1-x_0\|<\eps .
\]

\section*{Idea of the proof}

The key point is that no special topology of $C(K)$ is needed to find a
norming evaluation.  For any Banach space $E$, the injective tensor product
$E\ot Y$ embeds canonically into $C(B_{E^\ast},Y)$, where $B_{E^\ast}$ carries
the weak-star topology.  Hence every tensor element is a continuous
$Y$-valued function on a compact space and attains its norm somewhere.

Choose any $e\in S_E$ and embed $Y$ by $y\mapsto e\otimes y$.  If the BPBp in
$E\ot Y$ gives a norm-attaining approximant $\widetilde S$ close to
$e\otimes T$, evaluate $\widetilde Sx_1$ at a norming functional
$e_1^\ast\in B_{E^\ast}$.  The scalar $\alpha=e_1^\ast(e)$ must be close to
the unit circle because $\widetilde Sx_1$ is close to
$e\otimes Tx_1$.  Evaluation at $e_1^\ast$, followed by a scalar rotation,
therefore yields an exactly norm-attaining $Y$-valued operator close to $T$.

\begin{theorem}
Let $E$ be a nonzero Banach space.  Let $X$ and $Y$ be Banach spaces.  If the
pair $(X,E\ot Y)$ has the BPBp, then the pair $(X,Y)$ has the BPBp.

More precisely, if $(X,E\ot Y)$ has the BPBp with function $\eta$, then
$(X,Y)$ has the BPBp with function
\[
 \eps\longmapsto \eta(\eps/2).
\]
\end{theorem}

\begin{proof}
Fix $e\in S_E$.  We use the canonical isometric representation
\[
 E\ot Y \hookrightarrow C(B_{E^\ast},Y),
\]
where $B_{E^\ast}$ is weak-star compact and an elementary tensor
$a\otimes y$ is represented by the function
\[
 e^\ast\longmapsto e^\ast(a)y .
\]
For finite tensors this is exactly the injective norm formula; the completed
tensor product is the corresponding closed subspace of $C(B_{E^\ast},Y)$.

Let $0<\eps<1$ and put $\rho=\eps/2$.  Suppose
$T\in S_{\mathcal L(X,Y)}$ and $x_0\in S_X$ satisfy
\[
 \|Tx_0\|>1-\eta(\rho).
\]
Define $J:Y\to E\ot Y$ by $J(y)=e\otimes y$, and set
$\widetilde T=J\circ T$.  Since $\|e\|=1$, $J$ is an isometry, so
$\widetilde T\in S_{\mathcal L(X,E\ot Y)}$ and
\[
 \|\widetilde T x_0\|=\|Tx_0\|>1-\eta(\rho).
\]
Applying the BPBp for $(X,E\ot Y)$ with tolerance $\rho$, we get
$\widetilde S\in\mathcal L(X,E\ot Y)$ and $x_1\in S_X$ such that
\[
 \|\widetilde S\|=\|\widetilde Sx_1\|=1,\qquad
 \|\widetilde S-\widetilde T\|<\rho,\qquad
 \|x_1-x_0\|<\rho .
\]
As an element of $C(B_{E^\ast},Y)$, the function $\widetilde Sx_1$ attains
its norm.  Choose $e_1^\ast\in B_{E^\ast}$ such that
\[
 \|(\widetilde Sx_1)(e_1^\ast)\|=1.
\]
Let $\alpha=e_1^\ast(e)$.  Since
$\widetilde T x_1(e_1^\ast)=\alpha Tx_1$, we have
\[
 1=\|(\widetilde Sx_1)(e_1^\ast)\|
 \leq \|(\widetilde Sx_1-\widetilde T x_1)(e_1^\ast)\|
      +|\alpha|\,\|Tx_1\|
 < \rho+|\alpha|.
\]
Thus $|\alpha|>1-\rho>0$.  Choose $\omega\in\K$ with $|\omega|=1$ and
$\omega\alpha=|\alpha|$.

Define
\[
 S=\omega\,(e_1^\ast\otimes \mathrm{Id}_Y)\circ \widetilde S
 \in\mathcal L(X,Y).
\]
Then $\|S\|\leq 1$, while
\[
 \|Sx_1\|=\|(\widetilde Sx_1)(e_1^\ast)\|=1.
\]
Hence $\|S\|=\|Sx_1\|=1$.  Finally,
\[
\begin{aligned}
\|S-T\|
&\leq
\|\omega(e_1^\ast\otimes \mathrm{Id}_Y)(\widetilde S-\widetilde T)\|
  +\|\omega(e_1^\ast\otimes \mathrm{Id}_Y)\widetilde T-T\|  \\
&\leq \rho+\|\omega\alpha T-T\|
 = \rho+(1-|\alpha|)
 <2\rho
 =\eps .
\end{aligned}
\]
Also $\|x_1-x_0\|<\rho<\eps$.  Therefore $(X,Y)$ has the BPBp with modulus
$\eps\mapsto \eta(\eps/2)$.
\end{proof}

\section*{Same-modulus refinement}

The proof above loses a factor of two because the norming functional
$e_1^\ast$ need only satisfy $|e_1^\ast(e)|>1-\rho$, not $e_1^\ast(e)=1$.
If $E$ has a norm-one element $e$ whose state face
\[
K_e=\{e^\ast\in B_{E^\ast}:e^\ast(e)=1\}
\]
norms $E$, then the earlier argument over $K_e$ gives the same modulus
$\eta(\eps)$.  This is the state-norming-unit condition, equivalent to the
usual spear-vector formulation~\cite{KMMPO}; it includes closed unital
subspaces of $C(K)$ as a special case.

\section*{Scope and limitations}

This is a full answer for the qualitative BPBp descent problem associated with
injective tensor factors: every nonzero $E$ works.  It does not prove that the
same original modulus $\eta(\eps)$ is preserved for every $E$.  The same
modulus is proved here only under the additional state-norming-unit condition.

The separate $L_1(\mu,X)$ domain question in the same source paragraph is not
addressed by this packet.

\section*{Packet history}

This full packet supersedes two intermediate packets now preserved in the
local \texttt{history/} folder.  The first proved the same-modulus result for
closed unital subspaces of \(C(K)\).  The second replaced that concrete
hypothesis by the abstract state-norming-unit, or spear-vector, condition.
Both are now historical review aids; the ledger and run registry point to the
present full packet.

\section*{Novelty check}

The run indexes contained no exact packet or attempt for arXiv:1305.6420 and
no entry for the all-nonzero injective tensor factor formulation.  A bounded
web search for phrases including ``BPBp'', ``injective tensor'', ``$C(K,Y)$'',
``unital uniform algebra'', and ``spear vector'' found the source paper, the
spear-vector monograph~\cite{KMMPO}, and a related later preprint by Grando and
Santos~\cite{GS}, which proves a $C_0(L,Y)$ descent theorem and related
generalized AHSP stability results.  The downloaded arXiv source of that later
preprint did not state the all-nonzero injective tensor theorem proved here.

\section*{Human-review recommendation}

The key verifier focus is the canonical isometric embedding
$E\ot Y\subset C(B_{E^\ast},Y)$ and the scalar-rotation estimate turning
evaluation at $e_1^\ast$ into a norm-attaining operator within $\eps$ of $T$.
The proof is otherwise a direct quantitative refinement of the source paper's
Proposition 2.8 argument.

\begin{thebibliography}{9}
\bibitem{ACKLM}
R. Aron, Y. S. Choi, S. K. Kim, H. J. Lee, and M. Martin,
\emph{The Bishop-Phelps-Bollobas version of Lindenstrauss properties A and B},
Trans. Amer. Math. Soc. 367 (2015), 6085--6101; arXiv:1305.6420.

\bibitem{GS}
T. Grando and E. R. Santos,
\emph{Stability results of the Bishop-Phelps-Bollobas property and the
generalized AHSP}, arXiv:2506.17762.

\bibitem{KMMPO}
V. Kadets, M. Martin, J. Meri, and A. Perez,
\emph{Spear operators between Banach spaces}, arXiv:1701.02977.
\end{thebibliography}

\end{document}
